%!TEX root = ../template.tex
%%%%%%%%%%%%%%%%%%%%%%%%%%%%%%%%%%%%%%%%%%%%%%%%%%%%%%%%%%%%%%%%%%%%
%% abstract-pt.tex
%% NOVA thesis document file
%%
%% Abstract in Portuguese
%%%%%%%%%%%%%%%%%%%%%%%%%%%%%%%%%%%%%%%%%%%%%%%%%%%%%%%%%%%%%%%%%%%%

\typeout{NT FILE abstract-pt.tex}%



Os dispositivos \textit{wearables} de desporto comuns dependem predominantemente de dados cinemáticos provenientes de Unidades de Medição Inercial (\glspl{IMU}), não contendo o contexto fisiológico muscular que ajuda
a classificar movimentos atléticos complexos e de alta intensidade. Para colmatar esta limitação, esta dissertação apresenta o desenho, implementação e avaliação de um ecossistema de \textit{wearables} 
multimodal concebido especificamente para a análise de basquetebol. A arquitetura de hardware desenvolvida funde a cinemática de \glspl{IMU} de 9 eixos com os dados de sensores de Eletromiografia (\gls{EMG}), 
sendo transmitidos através de microcontroladores Arduino \gls{BLE} à frequência de 50Hz.

Para processar estes dados de alta frequência, foi desenvolvida uma aplicação Android dedicada para atuar como um centro descentralizado de computação. O software sincroniza os fluxos de dados de ambos os sensores 
numa janela deslizante de 300ms alimentando dois modelos de \gls{ML}: um classificador \textit{Random Forest} (\gls{RF}) sedeado num servidor para uma análise mais aprofundada, e uma rede neuronal TinyML quantizada, 
executada localmente para inferência de latência ultrabaixa.

A avaliação validou matematicamente a hipótese central da fusão de sensores. A análise de importância das características revelou que a \gls{VAR} e o valor \gls{RMS} da \gls{EMG} superaram a cinemática 
padrão do acelerómetro como os principais contribuidores na distinção de movimentos semelhantes no basquetebol. Em última análise, a implementação bem-sucedida de \textit{Edge} \gls{ML} eliminou 
totalmente a latência induzida pela rede, permitindo alertas instantâneos e ativos para a prevenção de cãimbras. Esta investigação demonstra que a harmonização entre sensores mecânicos, fisiologia elétrica e 
\gls{ML} incorporada constitui uma abordagem eficaz na salvaguarda da saúde de atletas.

% Palavras-chave
\keywords{
  \textit{Wearables} Desportivos \and
  Fusão de Sensores \and
  Eletromiografia (EMG) \and
  Unidade de Medição Inercial (IMU) \and
  \textit{Edge} ML \and
  TinyML \and
  Aprendizagem Automática (ML)
}